\documentclass[11pt,a4paper]{article}

% Packages
\usepackage[utf8]{inputenc}
\usepackage[T1]{fontenc}
\usepackage{amsmath,amssymb,amsfonts}
\usepackage{graphicx}
\usepackage{booktabs}
\usepackage{multirow}
\usepackage{hyperref}
\usepackage{cleveref}
\usepackage{algorithm}
\usepackage{algorithmic}
\usepackage[margin=1in]{geometry}
\usepackage{caption}
\usepackage{subcaption}
\usepackage{xcolor}

% Custom commands
\newcommand{\pcal}{p_{\text{cal}}}
\newcommand{\praw}{p_{\text{raw}}}
\newcommand{\Neff}{N_{\text{eff}}}
\newcommand{\abs}[1]{\left|#1\right|}

\title{Calibrated Large-Move Probability Estimation\\via Walk-Forward Monte Carlo with Online Recalibration}

\author{
  Chau Le \\
  \textit{Independent Research}
}

\date{\today}

\begin{document}

\maketitle

% =========================================================================
% ABSTRACT
% =========================================================================
\begin{abstract}
We present a framework for estimating the probability of large two-sided price moves (``Will the price move more than $\theta$\% in the next $H$ days?'') with calibrated uncertainty. The system couples GJR-GARCH volatility forecasts with Monte Carlo simulation (Student-$t$ innovations, Merton jump-diffusion, regime-conditional dynamics) and a three-layer online recalibrator (Platt scaling, multi-feature logistic regression, histogram post-calibration). A walk-forward protocol with explicit resolution queues enforces causal ordering. On four U.S.\ equities (SPY, GOOGL, AMZN, NVDA, 1997--2025) across horizons $H \in \{5, 10, 20\}$ evaluated by 5-fold expanding-window cross-validation, the framework passes all three promotion gates (ECE $\leq 0.02$, BSS $> 0$, AUC $\geq 0.55$) at the primary horizons $H\!=\!5$ and $H\!=\!10$ for three of four tickers (6 of 8 primary tests), with FDR-corrected paired block bootstrap tests rejecting the null of equal Brier loss against seven baselines at $p_{\text{adj}} < 0.001$ for SPY and most comparisons across the matrix. Hyperparameter search uses constrained Bayesian optimization (Optuna TPE with a constraints function and MedianPruner), cutting typical BO wall-clock time by 2--4$\times$. An economic significance analysis shows a simple probability-informed exposure rule reduces SPY maximum drawdown by 7--13\% relative while retaining 82--86\% of buy-and-hold returns, with the benefit surviving 5--10~bps transaction costs. All results are reproducible from the accompanying code, with raw forecasts timestamped to an append-only ledger for live out-of-sample verification.
\end{abstract}

\noindent\textbf{Keywords:} probability calibration, GARCH, Monte Carlo simulation, walk-forward validation, constrained Bayesian optimization, large-move prediction.

% =========================================================================
% 1. INTRODUCTION
% =========================================================================
\section{Introduction}
\label{sec:intro}

Position sizing, hedging, and entry timing all turn on a single quantity: the probability of a large near-term move. A trader sizing a \$1M exposure in a volatile stock wants to know the probability the price will move more than 5\% in the next two weeks, not simply the expected return or the 5\% VaR. If that probability is \emph{calibrated}---so that a predicted 15\% event happens approximately 15\% of the time---then rational decisions follow directly.

Existing approaches fall short in specific ways:
Value-at-Risk targets one-sided tail quantiles rather than two-sided move probabilities;
options-implied volatility embeds the variance risk premium and systematically overstates physical probabilities \cite{bollerslev2009expected};
GARCH-family models yield conditional variances but rarely close the loop to calibrated event probabilities;
and machine-learning classifiers often sacrifice calibration for discrimination \cite{guo2017calibration}.

We address this gap with an end-to-end framework whose contributions are:
\begin{enumerate}
  \item A \emph{calibrated probability} pipeline that achieves ECE $< 0.02$ at horizons $H\!=\!5$ and $H\!=\!10$ for three of four tested tickers, with honest analysis of where and why longer horizons degrade.
  \item A \emph{three-layer online recalibration} stack---Platt scaling, regime-conditional multi-feature logistic regression, histogram post-calibration---with systematic ablation showing each layer contributes measurably.
  \item A \emph{constrained Bayesian optimization} formulation of hyperparameter tuning: minimize ECE subject to the promotion-gate constraints (BSS $\geq 0$, AUC $\geq 0.55$), solved by TPE with a constraints function and a MedianPruner on the per-fold CV trajectory. This typically cuts BO wall-clock time by 2--4$\times$.
  \item An \emph{anti-overfitting discipline}: residual-based $\Neff$ estimation, Benjamini--Hochberg FDR correction across tests, block bootstrap with block size equal to the horizon, and five explicit overfitting diagnostics.
  \item A \emph{live verification} path: state-checkpointed prediction engine, asynchronous resolution, and an append-only forecast ledger with SHA-256 anchoring so the live out-of-sample track record is independently auditable.
\end{enumerate}

% =========================================================================
% 2. RELATED WORK
% =========================================================================
\section{Related Work}
\label{sec:related}

\paragraph{Probability calibration.}
\cite{platt1999probabilistic} introduced sigmoid-based post-hoc calibration; \cite{niculescu2005predicting} systematized comparisons with isotonic regression and histogram binning; \cite{guo2017calibration} showed modern classifiers are poorly calibrated despite high accuracy. Proper scoring rules \cite{gneiting2007strictly} decompose the Brier score into reliability, resolution, and uncertainty, and \cite{brocker2007increasing} showed reliability and sharpness jointly determine forecast quality. Our stack extends online Platt scaling to a multi-feature, regime-conditional setting designed for non-stationary financial time series where batch methods would require lookahead.

\paragraph{Volatility modeling and Monte Carlo.}
GARCH(1,1) \cite{bollerslev1986generalized} and its leverage-asymmetric GJR variant \cite{glosten1993relation} remain workhorses despite more complex alternatives \cite{hansen2005forecast}. Monte Carlo path simulation dates to \cite{boyle1977options}; Merton jump-diffusion \cite{merton1976option} adds discontinuous jumps; Student-$t$ innovations \cite{bollerslev1987conditionally} provide heavier tails. We use MC simulation as a structured probability-generation layer feeding the calibration stack, combining the inductive bias of GARCH-in-simulation dynamics with disciplined online recalibration.

\paragraph{Walk-forward evaluation and binary event prediction.}
\cite{tashman2000out} established walk-forward (rolling origin) as the standard for time-series forecast validation; \cite{bergmeir2012use} analyzed CV for time series and favored expanding-window protocols. For tail risk, \cite{christoffersen1998evaluating} proposed tests for VaR exceedances, \cite{engle2004caviar} introduced CAViaR, and \cite{patton2015good} developed tests for optimal VaR and ES forecasts. Unlike these one-sided tail-risk approaches, our target is two-sided move magnitude---a different prediction task relevant to position sizing and hedging.

% =========================================================================
% 3. METHODOLOGY
% =========================================================================
\section{Methodology}
\label{sec:methodology}

\subsection{Problem formulation}

Given daily closes $\{P_t\}_{t=1}^{T}$, for each date $t$ and horizon $H \in \{5, 10, 20\}$ we estimate
\begin{equation}
  \pcal(t, H) = \Pr\!\left(\abs{\frac{P_{t+H} - P_t}{P_t}} \geq \theta_H \,\middle|\, \mathcal{F}_t\right),
  \label{eq:target}
\end{equation}
where $\theta_H$ is a percentage threshold and $\mathcal{F}_t$ is the information set at time $t$. The binary label is $y_t = \mathbf{1}[\abs{r_{t:t+H}} \geq \theta_H]$, and calibration means $\Pr(y_t = 1 \mid \pcal = p) \approx p$ for all $p \in [0,1]$. Thresholds are ticker-specific, set by constrained BO (Section~\ref{sec:bo}). Baselines are evaluated against the full model's target labels to ensure identical event definitions.

\subsection{Volatility, simulation, and raw probabilities}

We fit GJR-GARCH(1,1) on returns:
\begin{equation}
  \sigma_{t+1}^2 = \omega + (\alpha + \gamma \mathbf{1}_{\epsilon_t < 0})\epsilon_t^2 + \beta\sigma_t^2,
\end{equation}
projecting to target persistence $\phi^\ast = 0.98$ whenever fitted $\phi = \alpha + \beta + \gamma/2 \geq 1$ (scaling $\alpha, \beta, \gamma$ by $s = \phi^\ast / \phi$ and setting $\omega = \sigma_t^2(1-\phi^\ast)$). On optimizer failure we fall back to EWMA volatility with a 252-day span \cite{jp1996riskmetrics}.

We simulate $N_\text{paths} = 30{,}000$ (boosted to 60{,}000 near the threshold) price paths in log-space with per-step GARCH dynamics:
\begin{align}
  \log P_{t+\Delta t} &= \log P_t + (\mu - \tfrac{1}{2}\sigma_t^2)\Delta t + \sigma_t\sqrt{\Delta t}\,Z_t + J_t, \\
  \sigma_{t+\Delta t}^2 &= \omega + (\alpha + \gamma\mathbf{1}_{Z_t<0})\sigma_t^2 Z_t^2 + \beta\sigma_t^2,
\end{align}
with $\Delta t = 1/252$, $Z_t \sim t_\nu / \sqrt{(\nu-2)/\nu}$ (unit-variance Student-$t$), and $J_t = \sum_{i=1}^{N_t} Y_i$, $N_t \sim \text{Pois}(\lambda\Delta t)$, $Y_i \sim \mathcal{N}(\mu_J, \sigma_J^2)$. Degrees of freedom $\nu$ vary by volatility regime ($\nu_\text{low}{=}10$, $\nu_\text{mid}{=}5$, $\nu_\text{high}{=}4$). The raw probability is the fraction of paths exceeding the threshold:
\begin{equation}
  \praw(t, H) = \frac{1}{N_\text{paths}} \sum_{i=1}^{N_\text{paths}} \mathbf{1}\!\left[\abs{\frac{P_{t+H}^{(i)} - P_t}{P_t}} \geq \theta_H\right].
\end{equation}

\subsection{Three-layer online recalibration}
\label{sec:calibration}

Raw MC probabilities are systematically biased by model misspecification and finite-sample estimation error. We apply three sequential calibration layers:

\textbf{(1) Online Platt scaling.} Logistic $\pcal = \sigma(a + b \cdot \text{logit}(\praw))$ with SGD updates $a \leftarrow a + \eta(y-\pcal)$, $b \leftarrow b + \eta(y-\pcal)\text{logit}(\praw)$, learning rate $\eta_n = \eta_0/\sqrt{1+n}$, 50-update warmup.

\textbf{(2) Multi-feature logistic regression.} Online logistic with $K \in \{6, 7, 8\}$ features and L2 regularization:
\begin{equation}
  \pcal = \sigma\!\left(\sum_{k=1}^{K} w_k x_k\right), \quad x = [\text{logit}(\praw), \sigma_t, \Delta\sigma_{20d}, \text{vol\_ratio}, \text{vol\_of\_vol}, 1, \text{earn?}, \text{iv\_ratio?}].
\end{equation}
Earnings proximity is active only at $H \leq 5$ for single stocks; implied-vol ratio is active when options data is available. Regime-conditional variants maintain separate weights per volatility regime.

\textbf{(3) Histogram post-calibration with Bayesian shrinkage.} Partition $[0, 1]$ into $B$ equal-width bins and apply additive corrections
\begin{equation}
  \delta_b = \frac{n_b}{n_b + \tau}(\bar{p}_b - \bar{y}_b), \qquad \pcal^{\text{post}} = \pcal - \delta_b,
  \label{eq:histogram}
\end{equation}
with Pool Adjacent Violators enforcing monotonicity. Bin count and prior strength $\tau$ are horizon-specific (more bins, weaker prior at $H\!=\!5$; fewer bins, stronger prior at $H\!=\!20$).

\subsection{Walk-forward protocol}

\Cref{alg:walkforward} shows the backtest engine. The resolution queue is the key no-lookahead mechanism: a prediction queued at time $t$ for horizon $H$ is resolved only when $t+H$ arrives, and calibrator updates consume only resolved outcomes.

\begin{algorithm}[htbp]
\caption{Walk-forward backtest with resolution queue}
\label{alg:walkforward}
\begin{algorithmic}[1]
\REQUIRE Prices $\{P_t\}$, config $\mathcal{C}$, horizons $\{H_1, \ldots, H_m\}$
\STATE Initialize calibrators $\mathcal{K}_H$ for each $H$; queue $\mathcal{Q} \leftarrow \emptyset$
\FOR{$t = t_\text{start}$ \textbf{to} $T$}
  \STATE \textbf{Resolve:} for each $(t', H, \praw) \in \mathcal{Q}$ with $t' + H \leq t$: compute $y$, update $\mathcal{K}_H$ with $(\praw, y)$, remove from $\mathcal{Q}$.
  \STATE \textbf{Estimate:} fit GARCH on $\{r_s\}_{s \leq t}$; extract $\sigma_t$.
  \FOR{each horizon $H$}
    \STATE Simulate $N$ MC paths; compute $\praw(t, H)$.
    \STATE $\pcal(t, H) \leftarrow \mathcal{K}_H.\text{calibrate}(\praw)$; push $(t, H, \praw)$ to $\mathcal{Q}$.
  \ENDFOR
\ENDFOR
\end{algorithmic}
\end{algorithm}

\subsection{Evaluation metrics}
\label{sec:metrics}

\textbf{Brier skill score.} $\text{BSS} = 1 - \text{Brier}_\text{model} / \text{Brier}_\text{clim}$ with $\text{Brier}_\text{clim} = \bar{y}(1-\bar{y})$; BSS $> 0$ beats the climatological base rate.

\textbf{AUC.} Discrimination---the probability that a randomly chosen event day is ranked above a randomly chosen non-event day.

\textbf{Expected calibration error.} $\text{ECE} = \sum_b (n_b/N)|\bar{p}_b - \bar{y}_b|$ with equal-width bins. ECE $\leq 0.02$ is our promotion threshold.

\textbf{Effective sample size.} Bartlett-type correction $\Neff = N / (1 + 2\sum_{k=1}^{H-1}(1-k/H)\hat{\rho}_k)$ with $\hat{\rho}_k$ computed on \emph{prediction residuals} $r_t = \pcal(t) - y_t$, not binary labels---ACF on 0/1 labels overestimates $\Neff$ by 30--40\%.

\textbf{Significance testing.} Paired circular block bootstrap \cite{efron1993introduction} on per-observation Brier losses with block size equal to $H$ and 2000 resamples; 95\% BCa intervals for BSS, AUC, ECE. All $p$-values are Benjamini--Hochberg FDR-corrected within each table \cite{benjamini1995controlling}.

\subsection{Constrained Bayesian optimization}
\label{sec:bo}

Hyperparameter tuning is formulated as a constrained optimization problem:
\begin{equation}
  \begin{aligned}
    \min_{\phi \in \Phi} \quad & \overline{\text{ECE}}(\phi) \\
    \text{s.t.} \quad & \text{ECE}_H(\phi) \leq 0.02 \quad \forall H, \\
                      & \text{AUC}_H(\phi) \geq 0.55 \quad \forall H, \\
                      & \text{BSS}_H(\phi) \geq 0 \quad \forall H,
  \end{aligned}
  \label{eq:cbo}
\end{equation}
where $\phi$ ranges over tunable parameters (thresholds, GARCH persistence target, multi-feature learning rate and L2) and $\overline{\text{ECE}}(\phi)$ is the mean pooled out-of-fold ECE across horizons on 5-fold expanding-window CV. We solve (\ref{eq:cbo}) with Optuna's Tree-structured Parzen Estimator (TPE) \cite{bergstra2011algorithms} configured with a \texttt{constraints\_func} that returns the horizon-worst slack $c = (\overline{\text{ECE}} - 0.02,\; 0.55 - \min_H \text{AUC}_H,\; -\min_H \text{BSS}_H)$ from each trial. TPE then models the density of feasible ($c \leq 0$) trials separately and focuses proposals on the feasible region.

Two further mechanisms cut wall-clock time. First, a \emph{MedianPruner} \cite{golovin2017google} tracks the running mean ECE across the five CV folds; after five warmup trials, any trial whose fold-$k$ running ECE exceeds the median of completed trials at fold $k$ is pruned before the remaining folds run. Second, a \emph{fast-fail} check terminates fold 0 immediately when its ECE exceeds $0.10$ on any horizon: a config that catastrophically miscalibrates on the first fold cannot be rescued by later folds. Together these cut typical BO wall-clock time by 2--4$\times$ without changing best-trial quality. Study versioning incorporates the feature-flag hash \emph{and} the BO-formulation version, so prior (unconstrained) studies do not contaminate the constrained search.

\subsection{Anti-overfitting controls}
\label{sec:antioverfitting}

Financial datasets are small relative to the space of configurations, so we apply five explicit controls:
(i) \emph{lean BO}: 6 tunable parameters (reduced from 14) with the remainder fixed at defaults;
(ii) \emph{adaptive event-rate guard} $r_\text{min} = \max(100 n_\text{params}/(2 N_\text{oof}),\; 0.03)$, ensuring $\geq 100$ events per parameter \cite{vittinghoff2007relaxing};
(iii) \emph{$\Neff$/parameter ratio} with GREEN/YELLOW/RED thresholds at $100\times / 50\times$;
(iv) \emph{generalization-gap monitoring} $\Delta = (\text{ECE}_\text{full} - \text{ECE}_\text{BO}) / \text{ECE}_\text{BO}$, GREEN if $\Delta < 0.25$; and
(v) \emph{temporal stability}: Brier scores in later CV folds are not systematically worse than earlier ones.

% =========================================================================
% 4. EXPERIMENTAL SETUP
% =========================================================================
\section{Experimental Setup}
\label{sec:experiments}

\subsection{Data and horizons}

We use daily closes for four U.S.\ equities spanning diverse capitalizations, sectors, and volatility regimes (\Cref{tab:data}). All series extend from the earliest available date (IPO or inception) through December 2025.

\begin{table}[htbp]
\centering
\small
\caption{Dataset summary. Ann.\ vol is annualized historical volatility.}
\label{tab:data}
\begin{tabular}{llrrr}
\toprule
Ticker & Period & Rows & Ann.\ vol & Sector \\
\midrule
SPY   & 2000--2025 & 6,538 & $\sim$16\% & Market ETF \\
GOOGL & 2004--2025 & 5,376 & $\sim$28\% & Technology \\
AMZN  & 1997--2025 & 7,202 & $\sim$55\% & Consumer/Tech \\
NVDA  & 1999--2025 & 6,777 & $\sim$60\% & Semiconductors \\
\bottomrule
\end{tabular}
\end{table}

Three horizons are evaluated: $H \in \{5, 10, 20\}$ trading days ($\approx 1, 2, 4$ weeks). $H\!=\!5$ and $H\!=\!10$ are \emph{primary}; $H\!=\!20$ is \emph{exploratory}, included to characterize where calibration begins to degrade. Thresholds $\theta_H$ are set per ticker by constrained BO with target event rates of 5--15\%.

\subsection{Baselines}

We compare against seven baselines spanning statistical, machine-learning, and practitioner approaches:
(1)~\textbf{Historical frequency}: rolling 252-day empirical event rate;
(2)~\textbf{GARCH-CDF}: parametric probability under Student-$t$, no MC;
(3)~\textbf{Implied-vol Black-Scholes}: log-normal probability from VIX-implied vol;
(4)~\textbf{Feature logistic}: online logistic on the same volatility features, no MC;
(5)~\textbf{Gradient boosting}: walk-forward \texttt{HistGradientBoostingClassifier} \cite{friedman2001greedy} with Platt-scaled outputs, refit every 63 days;
(6)~\textbf{VIX threshold rule}: a practitioner ``reduce exposure when VIX exceeds its 80th percentile'' heuristic;
(7)~\textbf{Market-implied straddle}: log-normal probability under VIX-implied volatility---what the options market thinks the probability is.

\subsection{Evaluation protocol}

\emph{Cross-validation}: 5-fold expanding-window, out-of-fold predictions pooled across folds for governance.
\emph{Promotion gates}: ECE $\leq 0.02$, BSS $> 0$, AUC $\geq 0.55$---all three must pass simultaneously.
\emph{Temporal hold-out}: walk-forward results partitioned at cutoff 2019-12-31. This is an \emph{online-adaptive} hold-out, not a frozen-model test: after the cutoff, GARCH and calibrators continue to update causally as post-cutoff labels resolve (each label at $t$ is available only at $t+H$). Post-cutoff data span COVID, Fed tightening, and the AI rally.

% =========================================================================
% 5. RESULTS
% =========================================================================
\section{Results}
\label{sec:results}

\subsection{Main cross-validation results}

\Cref{tab:main} shows 5-fold expanding-window CV with pooled OOF evaluation, 95\% BCa block bootstrap CIs, and FDR-corrected $p$-values. $\Neff$ is residual-based.

\begin{table}[htbp]
\centering
\small
\caption{Main results (pooled OOF). Bold = gate pass. $^{***}p_{\text{adj}}<0.01$, $^{**}p_{\text{adj}}<0.05$, $^{*}p_{\text{adj}}<0.10$ (Benjamini--Hochberg across all 12 tests). Top panel: primary horizons; bottom: exploratory $H\!=\!20$.}
\label{tab:main}
\begin{tabular}{llcccccl}
\toprule
Ticker & H & $\theta_H$ & BSS & AUC & ECE & $\Neff$ & Gates \\
\midrule
\multicolumn{8}{l}{\textit{Primary horizons}} \\
\midrule
\multirow{2}{*}{SPY}   & 5  & 3.72\%  & +0.089$^{***}$ & \textbf{0.779} & \textbf{0.0088} & 2{,}178 & \textbf{PASS} \\
                       & 10 & 5.54\%  & +0.055         & \textbf{0.749} & \textbf{0.0092} & 1{,}093 & \textbf{PASS} \\
\multirow{2}{*}{GOOGL} & 5  & 5.22\%  & +0.022$^{*}$   & \textbf{0.646} & \textbf{0.0099} & 1{,}418 & \textbf{PASS} \\
                       & 10 & 9.24\%  & +0.021         & \textbf{0.661} & \textbf{0.0121} &    770  & \textbf{PASS} \\
\multirow{2}{*}{AMZN}  & 5  & 8.52\%  & +0.029         & \textbf{0.705} & \textbf{0.0076} & 1{,}881 & \textbf{PASS} \\
                       & 10 & 15.99\% & +0.020         & \textbf{0.700} & \textbf{0.0046} & 1{,}075 & \textbf{PASS} \\
\multirow{2}{*}{NVDA}  & 5  & 12.18\% & $-$0.001       & 0.604          & \textbf{0.0021} & 1{,}845 & FAIL \\
                       & 10 & 18.22\% & $-$0.011       & 0.588          & \textbf{0.0095} &    944  & FAIL \\
\midrule
\multicolumn{8}{l}{\textit{Exploratory horizon}} \\
\midrule
SPY    & 20 & 7.05\%  & +0.051         & \textbf{0.779} & 0.0265          & 632 & FAIL \\
GOOGL  & 20 & 12.65\% & +0.024         & \textbf{0.672} & \textbf{0.0191} & 422 & \textbf{PASS} \\
AMZN   & 20 & 20.50\% & +0.008         & \textbf{0.708} & \textbf{0.0132} & 606 & \textbf{PASS} \\
NVDA   & 20 & 25.66\% & $-$0.017       & 0.562          & 0.0264          & 418 & FAIL \\
\bottomrule
\end{tabular}
\end{table}

Three of four tickers (SPY, GOOGL, AMZN) pass all promotion gates at primary horizons---6 of 8 primary tests. SPY $H\!=\!5$ reaches FDR-corrected significance; other passing configurations have raw $p < 0.10$ but lose individual significance under conservative multi-test correction. AMZN produces the strongest calibration ($H\!=\!10$ ECE of 0.0046) despite the highest annualized volatility, reflecting the benefit of 28 years of data spanning diverse regimes. NVDA fails primary horizons with slightly negative BSS---the threshold-selection failure mode discussed in Section~\ref{sec:nvda}. At $H\!=\!20$ calibration degrades for SPY and NVDA but holds for GOOGL and AMZN; see Section~\ref{sec:horizon}.

\subsection{Baseline comparison}

\Cref{tab:baselines} reports SPY as representative. Across the full 4-ticker matrix, the full model is FDR-significant against most baselines; a few AMZN comparisons against the VIX rule and straddle baseline at $H\!=\!10$/$H\!=\!20$ do not reach FDR significance, which is reported honestly in the machine-readable tables.

\begin{table}[htbp]
\centering
\small
\caption{Full model vs.\ baselines (SPY, representative). $p$-values are paired block bootstrap (block size $= H$, 2000 resamples) with FDR correction. All methods evaluated on the full model's target labels.}
\label{tab:baselines}
\begin{tabular}{llcccc}
\toprule
Method & H & BSS & AUC & ECE & $p$-value \\
\midrule
Full Model             & 5  & +0.131 & 0.814 & 0.013 & --- \\
Historical Frequency   & 5  & $-$0.002 & 0.726 & 0.034 & $<$0.001$^{***}$ \\
GARCH-CDF              & 5  & $-$0.002 & 0.745 & 0.033 & $<$0.001$^{***}$ \\
Feature Logistic       & 5  & $-$0.563 & 0.629 & 0.188 & $<$0.001$^{***}$ \\
Gradient Boosting      & 5  & $-$0.067 & 0.723 & 0.058 & $<$0.001$^{***}$ \\
VIX Threshold Rule     & 5  & +0.049 & 0.749 & 0.014 & $<$0.001$^{***}$ \\
Market-Implied         & 5  & $-$0.121 & 0.815 & 0.097 & $<$0.001$^{***}$ \\
\midrule
Full Model             & 10 & +0.106 & 0.810 & 0.013 & --- \\
Historical Frequency   & 10 & $-$0.022 & 0.698 & 0.025 & $<$0.001$^{***}$ \\
GARCH-CDF              & 10 & $-$0.054 & 0.739 & 0.039 & $<$0.001$^{***}$ \\
Feature Logistic       & 10 & $-$0.787 & 0.600 & 0.193 & $<$0.001$^{***}$ \\
Gradient Boosting      & 10 & $-$0.081 & 0.714 & 0.054 & $<$0.001$^{***}$ \\
VIX Threshold Rule     & 10 & +0.027 & 0.714 & 0.009 & $<$0.001$^{***}$ \\
Market-Implied         & 10 & $-$0.222 & 0.807 & 0.095 & $<$0.001$^{***}$ \\
\bottomrule
\end{tabular}
\end{table}

Two comparisons merit emphasis. The practitioner \emph{VIX threshold rule} is the simplest viable baseline and achieves positive BSS, but the full model's BSS is significantly higher ($p_{\text{adj}} < 0.001$). The \emph{market-implied straddle} probability has the strongest discrimination (AUC 0.815 at $H\!=\!5$) but the worst calibration (ECE 0.097) of any method---as expected, since implied volatility embeds the variance risk premium. Our calibration stack extracts the informational content of implied vol while correcting its risk-neutral bias.

\subsection{Ablation}

\Cref{tab:ablation} (SPY $H\!=\!5$) shows the cumulative contribution of each pipeline component.

\begin{table}[htbp]
\centering
\small
\caption{Ablation study: cumulative component contributions (SPY $H\!=\!5$).}
\label{tab:ablation}
\begin{tabular}{lccc}
\toprule
Configuration & BSS & AUC & ECE \\
\midrule
A: Base GBM (constant vol, Platt)     & +0.063 & 0.747 & 0.012 \\
B: +GARCH-in-sim                      & +0.072 & 0.755 & 0.015 \\
C: +Student-$t$ fat tails             & +0.081 & 0.758 & 0.015 \\
D: +Multi-feature calibration         & +0.067 & 0.759 & 0.027 \\
E: +Histogram post-calibration        & +0.081 & 0.762 & 0.008 \\
F: +Implied volatility                & +0.090 & 0.779 & 0.009 \\
\bottomrule
\end{tabular}
\end{table}

\emph{Histogram post-calibration} produces the largest ECE improvement ($\Delta\text{ECE} = -0.019$): it absorbs the residual errors from the multi-feature logistic layer. Step D \emph{temporarily} degrades ECE (0.015 $\to$ 0.027) by adding features the histogram layer must then regularize. \emph{Implied volatility} produces the largest AUC improvement ($\Delta\text{AUC} = +0.017$): forward-looking vol from options complements realized-vol features. \emph{Student-$t$ tails} and \emph{GARCH-in-sim} each add $\Delta\text{BSS} \approx +0.009$.

\subsection{Temporal hold-out}

\Cref{tab:holdout} presents online-adaptive hold-out results (train through 2019-12-31, evaluate 2020--2025).

\begin{table}[htbp]
\centering
\small
\caption{Online-adaptive hold-out: train through 2019-12-31, evaluate 2020--2025. Bold ECE $\leq 0.02$.}
\label{tab:holdout}
\begin{tabular}{llcccc}
\toprule
Ticker & H & BSS & AUC & ECE & Era breakdown (ECE) \\
\midrule
SPY   & 5  & +0.119 & 0.776 & \textbf{0.008} & COVID 0.062, Tight 0.008, Post 0.021 \\
SPY   & 10 & +0.076 & 0.725 & \textbf{0.018} & COVID 0.063, Tight 0.031, Post 0.006 \\
GOOGL & 5  & +0.014 & 0.622 & \textbf{0.012} & COVID 0.078, Tight 0.019, Post 0.021 \\
AMZN  & 5  & +0.007 & 0.629 & \textbf{0.006} & COVID 0.016, Tight 0.014, Post 0.009 \\
\bottomrule
\end{tabular}
\end{table}

Pooled hold-out ECE stays below 0.02 for all four configurations shown. The era breakdown makes clear that COVID-2020 is the hardest sub-period (ECE up to 0.078), consistent with the unprecedented realized volatility and the limited 220-sample window. Post-2023 calibration is strongest, indicating that the online update mechanism adapts to regime shifts, albeit slowly during extreme transitions.

\subsection{Economic significance}

We evaluate a simple risk-managed strategy: reduce SPY exposure to 50\% when $\pcal(H\!=\!5) > 0.20$.

\begin{table}[htbp]
\centering
\small
\caption{Economic significance (SPY, $H\!=\!5$, $\pcal > 0.20$ trigger).}
\label{tab:economic}
\begin{tabular}{lcc}
\toprule
Metric                 & Buy-and-Hold & Risk-Managed \\
\midrule
Annualized Return      & 9.6\%        & 8.3\% \\
Annualized Volatility  & 19.3\%       & 16.1\% \\
Sharpe Ratio           & 0.50         & 0.52 \\
Max Drawdown           & $-$55.2\%    & $-$48.0\% \\
CVaR (95\%)            & 2.9\%        & 2.5\% \\
\% days at reduced weight & ---        & 10\% \\
\bottomrule
\end{tabular}
\end{table}

The strategy retains 86\% of buy-and-hold returns while reducing maximum drawdown by 7.2 percentage points (13\% relative) and modestly improving Sharpe. With realistic transaction costs, Sharpe degrades from 0.52 at 0~bps to 0.50 at 5~bps and 0.46 at 20~bps, with the drawdown-reduction benefit surviving to 10~bps (9\% reduction) but narrowing at 20~bps (6\% reduction). Turnover is low---position changes on roughly 10\% of days.

% =========================================================================
% 6. DISCUSSION
% =========================================================================
\section{Discussion}
\label{sec:discussion}

\subsection{Why longer horizons are harder}
\label{sec:horizon}

At $H\!=\!20$, three compounding forces degrade calibration:
(i) fewer non-overlapping observations reduce $\Neff$ to 40--60\% of the $H\!=\!5$ value;
(ii) the mapping from current volatility to $H$-day move probability grows more non-linear as more regime transitions fit in the prediction window; and
(iii) GARCH persistence pulls the conditional variance toward the unconditional level, shrinking the discriminative signal.
Horizon-specific histogram parameterization (fewer bins, stronger prior at $H\!=\!20$) partially compensates, but the finite-data limit remains. We therefore scope primary claims to $H\!=\!5$ and $H\!=\!10$ and present $H\!=\!20$ as an exploratory boundary.

\subsection{NVDA: threshold-selection pathology}
\label{sec:nvda}

NVDA fails all three horizons despite 6,777 daily observations and 30 BO trials. The root cause is threshold selection: BO-tuned thresholds of 12--26\% produce event rates of 4--6\%, and the model's BSS is slightly negative at every horizon. Hold-out validation confirms severe overfitting: train CV shows 3/3 PASS (mean ECE 0.006) but hold-out shows 0/3 PASS (mean ECE 0.044), a 648\% generalization gap. BO's anti-overfitting controls correctly flag the pathology but cannot resolve it without widening the threshold search range. We report NVDA honestly to characterize the framework's boundary conditions.

\subsection{From backtest to live prediction}

A common objection to financial backtesting work is that the systems cannot generate live predictions. We close this gap end-to-end:
(i)~\emph{state checkpointing}---at the end of a walk-forward run, all calibrator states (weight vectors, histogram corrections, learning rates), GARCH parameters, and metadata are serialized to JSON;
(ii)~a \emph{prediction engine} loads a checkpoint and generates forward-looking probabilities from new price data in seconds;
(iii)~\emph{asynchronous resolution} updates the log with outcomes at $t+H$ for continuous calibrator refinement;
(iv)~a \emph{daily runner} publishes forecasts to an append-only JSONL ledger with SHA-256 manifest anchoring and a static verification dashboard.

The live engine uses a simplified feature pipeline relative to the full backtest (no implied-vol blending or state-dependent jumps in the current implementation); extending it to full parity is future work.

\subsection{Limitations}

\begin{enumerate}
  \item \textbf{Overlapping windows.} $H$-step predictions from consecutive days share $H{-}1$ outcome days, creating autocorrelated labels. We track $\Neff$ via residual-based ACF but statistical power is reduced.
  \item \textbf{Cold start.} The calibrator requires $\sim$100 resolved outcomes before activating; early predictions are uncalibrated.
  \item \textbf{Horizon scope.} Primary claims are limited to $H \in \{5, 10\}$; $H\!=\!20$ degrades for SPY and NVDA.
  \item \textbf{Threshold dependence.} The binary event definition determines both labels and base rates; results are sensitive to $\theta$.
  \item \textbf{Online-adaptive hold-out.} GARCH and calibrators keep updating causally post-cutoff; this is not a frozen-model test.
  \item \textbf{Live engine parity.} The live engine omits implied-vol blending and state-dependent jumps; live numbers should not be assumed equal to backtest.
  \item \textbf{Direction-agnostic.} The system predicts magnitude, not direction.
  \item \textbf{Approximate transaction costs.} Proportional costs only; no market impact, variable spreads, or timing effects.
  \item \textbf{Implied-vol proxies.} Single-name ``market-implied'' baselines use VIX as a market-wide proxy, weakening the comparison for single stocks.
\end{enumerate}

Auxiliary robustness analyses---prediction sharpness, conditional ECE by era and volatility regime, universal-threshold sensitivity (5\% and 10\% fixed thresholds), cross-asset correlation of daily $\pcal$, and an out-of-sample AAPL check---are reported in Appendix~\ref{app:robustness}.

% =========================================================================
% 7. CONCLUSION
% =========================================================================
\section{Conclusion}
\label{sec:conclusion}

We presented an end-to-end framework for calibrated large-move probability estimation combining GARCH-driven Monte Carlo simulation with three-layer online recalibration and a constrained Bayesian optimization procedure for hyperparameter tuning. At primary horizons ($H\!=\!5$, $H\!=\!10$) the system achieves ECE $< 0.02$ for three of four tickers with FDR-significant Brier improvements over seven baselines for the representative SPY comparison. Ablation confirms that each pipeline layer contributes measurably; temporal hold-out shows useful calibration through COVID, Fed tightening, and the AI rally with post-2023 being strongest; and a simple probability-informed exposure rule produces economically meaningful drawdown reduction that survives realistic transaction costs. The live prediction engine---with state checkpointing, asynchronous resolution, and a tamper-evident append-only ledger---closes the backtest-to-production loop for independent verification.

The core insight is that \emph{generating} probabilities (via MC simulation) and \emph{calibrating} them (via online learning) are complementary: simulation supplies the structured inductive bias about volatility dynamics, while calibration ensures that signal translates into trustworthy probabilities. Future work: extending density calibration beyond binary events, incorporating directional signals, ticker-specific implied-volatility surfaces, and cross-asset joint modeling for portfolio-level risk.

Reproduction code, data snapshots, and environment specifications are available at \url{https://github.com/chaulemuoichin/calibrated-large-move-probability-engine}.

% =========================================================================
% APPENDIX
% =========================================================================
\appendix

\section{Additional Robustness Analyses}
\label{app:robustness}

\paragraph{Sharpness.}
A model that always predicts near the base rate achieves low ECE trivially. For SPY $H\!=\!5$ (base rate $\sim$5\%), the model's $\pcal$ ranges 0.00--0.68 with standard deviation $\sim$0.10 and $\sim$52\% of predictions deviating more than 5~percentage points from the base rate---confirming meaningful discrimination rather than base-rate mimicry.

\paragraph{Conditional ECE.}
We compute ECE conditioned on volatility terciles (rolling 20-day realized vol) and on time eras (pre-2020, COVID-2020, tightening 2021--22, post-2023). Calibration is worst during COVID-2020 and in the high-vol tercile, as expected, but even in the worst bucket SPY $H\!=\!5$ ECE stays below 0.04. Post-2023 is consistently the best, indicating that online learning adapts to the current regime.

\paragraph{Universal-threshold sensitivity.}
To test generalizability beyond BO-tuned thresholds, we evaluate fixed universal thresholds of 5\% and 10\% across all four tickers. At 5\%: SPY passes $H\!=\!5$ and $H\!=\!10$; GOOGL passes $H\!=\!5$; AMZN and NVDA fail (event rates too high to calibrate). At 10\%: SPY and GOOGL pass $H\!=\!5$ and $H\!=\!10$; NVDA passes $H\!=\!5$; AMZN still fails. This confirms calibration skill without per-ticker tuning (5/24 universal-threshold tests pass) while showing that ticker-specific thresholds improve results substantially.

\paragraph{Cross-asset correlation.}
Pairwise Pearson correlation of daily $\pcal$ across tickers at $H\!=\!5$ ranges 0.37 (SPY--AMZN) to 0.43 (SPY--NVDA), while within-ticker cross-horizon correlations exceed 0.90---both driven by shared exposure to VIX-level vol shocks. A portfolio using these predictions must account for this dependence; joint modeling is left to future work.

\paragraph{AAPL out-of-sample check.}
AAPL (6,538 daily observations, not in the primary ticker set) passes all three gates at all three horizons ($H\!=\!5$: ECE 0.014, BSS +0.075, AUC 0.742; $H\!=\!10$: ECE 0.013, BSS +0.059, AUC 0.769; $H\!=\!20$: ECE 0.019, BSS +0.061, AUC 0.756), suggesting the framework generalizes to additional tickers when thresholds produce adequate event rates.

% =========================================================================
% REFERENCES
% =========================================================================
\begin{thebibliography}{25}

\bibitem{benjamini1995controlling}
Y.~Benjamini and Y.~Hochberg.
\newblock Controlling the false discovery rate: A practical and powerful approach to multiple testing.
\newblock \emph{Journal of the Royal Statistical Society: Series B}, 57(1):289--300, 1995.

\bibitem{bergmeir2012use}
C.~Bergmeir and J.~M. Ben\'{\i}tez.
\newblock On the use of cross-validation for time series predictor evaluation.
\newblock \emph{Information Sciences}, 191:192--213, 2012.

\bibitem{bergstra2011algorithms}
J.~Bergstra, R.~Bardenet, Y.~Bengio, and B.~K\'{e}gl.
\newblock Algorithms for hyper-parameter optimization.
\newblock In \emph{Advances in Neural Information Processing Systems}, volume~24, 2011.

\bibitem{bollerslev1986generalized}
T.~Bollerslev.
\newblock Generalized autoregressive conditional heteroskedasticity.
\newblock \emph{Journal of Econometrics}, 31(3):307--327, 1986.

\bibitem{bollerslev1987conditionally}
T.~Bollerslev.
\newblock A conditionally heteroskedastic time series model for speculative prices and rates of return.
\newblock \emph{The Review of Economics and Statistics}, 69(3):542--547, 1987.

\bibitem{bollerslev2009expected}
T.~Bollerslev, G.~Tauchen, and H.~Zhou.
\newblock Expected stock returns and variance risk premia.
\newblock \emph{The Review of Financial Studies}, 22(11):4463--4492, 2009.

\bibitem{boyle1977options}
P.~P. Boyle.
\newblock Options: A {Monte Carlo} approach.
\newblock \emph{Journal of Financial Economics}, 4(3):323--338, 1977.

\bibitem{brocker2007increasing}
J.~Br{\"o}cker and L.~A. Smith.
\newblock Increasing the reliability of reliability diagrams.
\newblock \emph{Weather and Forecasting}, 22(3):651--661, 2007.

\bibitem{christoffersen1998evaluating}
P.~F. Christoffersen.
\newblock Evaluating interval forecasts.
\newblock \emph{International Economic Review}, 39(4):841--862, 1998.

\bibitem{efron1993introduction}
B.~Efron and R.~J. Tibshirani.
\newblock \emph{An Introduction to the Bootstrap}.
\newblock Chapman and Hall/CRC, 1993.

\bibitem{friedman2001greedy}
J.~H. Friedman.
\newblock Greedy function approximation: A gradient boosting machine.
\newblock \emph{The Annals of Statistics}, 29(5):1189--1232, 2001.

\bibitem{engle2004caviar}
R.~F. Engle and S.~Manganelli.
\newblock {CAViaR}: Conditional autoregressive value at risk by regression quantiles.
\newblock \emph{Journal of Business \& Economic Statistics}, 22(4):367--381, 2004.

\bibitem{glosten1993relation}
L.~R. Glosten, R.~Jagannathan, and D.~E. Runkle.
\newblock On the relation between the expected value and the volatility of the nominal excess return on stocks.
\newblock \emph{The Journal of Finance}, 48(5):1779--1801, 1993.

\bibitem{gneiting2007strictly}
T.~Gneiting and A.~E. Raftery.
\newblock Strictly proper scoring rules, prediction, and estimation.
\newblock \emph{Journal of the American Statistical Association}, 102(477):359--378, 2007.

\bibitem{golovin2017google}
D.~Golovin, B.~Solnik, S.~Moitra, G.~Kochanski, J.~Karro, and D.~Sculley.
\newblock {Google Vizier}: A service for black-box optimization.
\newblock In \emph{Proceedings of the 23rd ACM SIGKDD International Conference on Knowledge Discovery and Data Mining}, pages 1487--1495, 2017.

\bibitem{guo2017calibration}
C.~Guo, G.~Pleiss, Y.~Sun, and K.~Q. Weinberger.
\newblock On calibration of modern neural networks.
\newblock In \emph{Proceedings of the 34th International Conference on Machine Learning}, pages 1321--1330, 2017.

\bibitem{hansen2005forecast}
P.~R. Hansen and A.~Lunde.
\newblock A forecast comparison of volatility models: does anything beat a {GARCH}(1,1)?
\newblock \emph{Journal of Applied Econometrics}, 20(7):873--889, 2005.

\bibitem{jp1996riskmetrics}
{J.P. Morgan}.
\newblock {RiskMetrics} -- Technical Document.
\newblock Technical report, J.P. Morgan/Reuters, 1996.

\bibitem{merton1976option}
R.~C. Merton.
\newblock Option pricing when underlying stock returns are discontinuous.
\newblock \emph{Journal of Financial Economics}, 3(1-2):125--144, 1976.

\bibitem{niculescu2005predicting}
A.~Niculescu-Mizil and R.~Caruana.
\newblock Predicting good probabilities with supervised learning.
\newblock In \emph{Proceedings of the 22nd International Conference on Machine Learning}, pages 625--632, 2005.

\bibitem{patton2015good}
A.~J. Patton.
\newblock Evaluating volatility and correlation forecasts.
\newblock In \emph{Handbook of Economic Forecasting}, volume~2, pages 801--838. Elsevier, 2015.

\bibitem{platt1999probabilistic}
J.~Platt.
\newblock Probabilistic outputs for support vector machines and comparisons to regularized likelihood methods.
\newblock In \emph{Advances in Large Margin Classifiers}, pages 61--74. MIT Press, 1999.

\bibitem{tashman2000out}
L.~J. Tashman.
\newblock Out-of-sample tests of forecasting accuracy: an analysis and review.
\newblock \emph{International Journal of Forecasting}, 16(4):437--450, 2000.

\bibitem{vittinghoff2007relaxing}
E.~Vittinghoff and C.~E. McCulloch.
\newblock Relaxing the rule of ten events per variable in logistic and {Cox} regression.
\newblock \emph{American Journal of Epidemiology}, 165(6):710--718, 2007.

\end{thebibliography}

\end{document}
